\medskip

\emph{Les parties A et B peuvent être traitées indépendamment}

\medskip
\textbf{Partie A}
\medskip

On considère la fonction $f$ définie sur $\R$ par

\[f(x) = \left(x + \dfrac12\right)\e^{-x} + x.\]

\medskip

\begin{enumerate}
\item Déterminer les limites de $f$ en $- \infty$ et en $+ \infty$.
\item On admet que $f$ est deux fois dérivable sur $\R$.
	\begin{enumerate}
		\item Démontrer que, pour tout $x \in \R$,
		
\[f''(x) = \left(x - \dfrac32\right) \e^{-x}.\]

		\item En déduire les variations et le minimum de la fonction $f'$ sur $\R$.
		\item Justifier que pour tout $x \in \R,\: f'(x) > 0$.
		\item En déduire que l'équation $f(x) = 0$ admet une unique solution sur $\R$.
		\item Donner une valeur arrondie à $10^{-3}$ de cette solution.
	\end{enumerate}
\end{enumerate}

\medskip
\textbf{Partie B}
\medskip

On considère une fonction $h$, définie et dérivable sur $\R$, ayant une expression de la forme 

\begin{center}$h(x)= (ax + b)\e^{-x} + x$, \:où $a$ et $b$\: sont deux réels.\end{center}

Dans un repère orthonormé ci-après figurent:

\begin{itemize}
\item[$\bullet$~~]la courbe représentative $\mathcal{C}_h$ de la fonction $h$ ;
\item[$\bullet$~~]les points A et B de coordonnées respectives $(-2~;~-2,5)$ et (2~;~3,5).
\end{itemize}

\begin{center}
\psset{unit=0.8cm}
\begin{pspicture*}(-5.9,-4.9)(5.9,5.9)
\psaxes[linewidth=1.25pt,labelFontSize=\scriptstyle](0,0)(-5.9,-4.9)(5.9,5.9)
\psplot[plotpoints=500,linewidth=1.25pt]{-4}{4}{1.5 x mul 0.5 add}
\psplot[plotpoints=2000,linewidth=1.25pt,linecolor=cyan]{-4}{4}{x 0.5 add 2.71828 x exp div x add}
\psdots(-2,-2.5)(2,3.5)
\uput[r](-2,-2.5){A}\uput[r](2,3.5){B}\uput[dr](3.8,4){\cyan $\mathcal{C}_h$}
\end{pspicture*}
\end{center}

\begin{enumerate}
\item Conjecturer, avec la précision permise par le graphique, les abscisses des éventuels points d'inflexion de la courbe représentative de la fonction $h$.
\item Sachant que la fonction $h$ admet sur $\R$ une dérivée seconde d'expression

\[h''(x) = - \dfrac32 \e^{- x} + x\e^{- x}.\]

valider ou non la conjecture précédente.
\item Déterminer une équation de la droite (AB).
\item Sachant que la droite (AB) est tangente à la courbe représentative de la fonction
$h$ au point d'abscisse 0, en déduire les valeurs de $a$ et $b$.
\end{enumerate}

\bigskip
